\documentclass[11pt]{article}

\usepackage[T1]{fontenc}
\usepackage[utf8]{inputenc}
\usepackage{lmodern}
\usepackage[margin=1in]{geometry}
\usepackage{microtype}
\usepackage{amsmath,amssymb,amsthm,mathtools}
\usepackage{booktabs}
\usepackage{enumitem}
\usepackage{longtable}
\usepackage{tabularx}
\usepackage{xcolor}
\usepackage{hyperref}
\usepackage[nameinlink,noabbrev]{cleveref}

\definecolor{linkblue}{HTML}{174A7E}
\hypersetup{
  colorlinks=true,
  linkcolor=linkblue,
  citecolor=linkblue,
  urlcolor=linkblue,
  pdftitle={A Detailed Proof of the McKay Equality},
  pdfauthor={Clawristotle contributors}
}

\newtheorem{theorem}{Theorem}[section]
\newtheorem{proposition}[theorem]{Proposition}
\newtheorem{lemma}[theorem]{Lemma}
\newtheorem{definition}[theorem]{Definition}
\newtheorem{remark}[theorem]{Remark}

\newcommand{\Irr}{\operatorname{Irr}}
\newcommand{\Aut}{\operatorname{Aut}}
\newcommand{\Inn}{\operatorname{Inn}}
\newcommand{\Out}{\operatorname{Out}}
\newcommand{\Norm}{\operatorname{N}}
\newcommand{\Cent}{\operatorname{C}}
\newcommand{\Syl}{\operatorname{Syl}}
\newcommand{\Z}{\operatorname{Z}}
\newcommand{\Ind}{\operatorname{Ind}}
\newcommand{\Res}{\operatorname{Res}}
\newcommand{\IMK}{\mathrm{iMK}}
\newcommand{\Ad}{\mathrm{A}(d)}
\newcommand{\Bd}{\mathrm{B}(d)}
\newcommand{\Ainf}{\mathrm{A}(\infty)}
\newcommand{\F}{\mathbf{F}}
\newcommand{\G}{\mathbf{G}}
\newcommand{\M}{\mathbf{M}}
\newcommand{\Sbf}{\mathbf{S}}

\title{\textbf{A Detailed Natural-Language Proof of the McKay Equality}\\
\large A proof certificate through the inductive McKay condition}
\author{Clawristotle contributors}
\date{July 2026}

\begin{document}

\maketitle

\begin{abstract}
For a finite group \(X\), a prime \(\ell\), and a Sylow
\(\ell\)-subgroup \(S\), the McKay equality says that \(X\) and
\(\Norm_X(S)\) have the same number of irreducible complex characters
whose degrees are not divisible by \(\ell\).  The first complete proof is
not a short argument: it combines the Isaacs--Malle--Navarro reduction,
the classification of finite simple groups, a family-by-family verification
of the inductive McKay condition, and the final type-\(D\) analysis of
Cabanes and Sp{\"a}th.

This document gives a detailed proof at the correct level of granularity:
every major imported theorem is stated, every family in the classification
is accounted for, the final type-\(D\) trichotomy is expanded, and the last
induction from an intermediate local subgroup to the Sylow normalizer is
proved explicitly.  It is a proof by fully identified literature results,
not a claim to reproduce the several hundred pages of character theory and
the classification on which those results depend.
\end{abstract}

{\small\tableofcontents}
\clearpage

\section{Statement and conventions}

Let \(X\) be a finite group.  An ordinary complex character of \(X\) is
the trace of a finite-dimensional complex representation.  Write
\(\Irr(X)\) for the set of characters afforded by irreducible
representations, with isomorphic representations counted once.  The
degree of \(\chi\in\Irr(X)\) is
\[
  \chi(1)=\dim_{\mathbb C} V_\chi\in\mathbb Z_{>0}.
\]
For a prime \(\ell\), define
\[
  \Irr_{\ell'}(X)
  :=\{\chi\in\Irr(X)\mid \ell\nmid \chi(1)\}.
\]

\begin{theorem}[McKay equality]\label{thm:mckay}
Let \(X\) be a finite group, let \(\ell\) be prime, and let
\(S\in\Syl_\ell(X)\).  Then
\[
  \bigl|\Irr_{\ell'}(X)\bigr|
  =
  \bigl|\Irr_{\ell'}(\Norm_X(S))\bigr|.
\]
\end{theorem}

The theorem does not depend on the chosen Sylow subgroup: any two Sylow
\(\ell\)-subgroups are conjugate, conjugation identifies their
normalizers, and a group isomorphism transports irreducible characters
without changing their degrees.

\section{The stronger inductive condition}

\subsection{Character triples}

A \emph{character triple} is a triple \((A,X,\chi)\) in which
\(X\trianglelefteq A\) and \(\chi\in\Irr(X)\) is invariant under \(A\).
Given a second triple \((H,M,\chi')\) with
\[
  \Cent_A(X)\leq H\leq A,\qquad A=XH,\qquad H\cap X=M,
\]
the relation
\[
  (A,X,\chi)\geq_c(H,M,\chi')
\]
of centrally isomorphic character triples records substantially more than
an equality of cardinalities.  Definition 2.5 of
\cite{CabanesSpath2026} requires associated projective representations
of \(A\) and \(H\) whose factor sets agree on \(H\times H\), and whose
scalar matrices agree on \(\Cent_A(X)\).  The relation supplies character
correspondences over the normal subgroups and has three structural
properties used below:

\begin{enumerate}[label=(\roman*)]
  \item it is transitive;
  \item it can be restricted to suitable subgroups of the outer action;
  \item its two sides can be transported through the ``Butterfly''
        construction when they induce the same automorphisms.
\end{enumerate}

The precise relation and these closure properties are developed in
Navarro--Sp{\"a}th's language and in Rossi's reformulation
\cite{Rossi2023,Navarro2018}.  They are the mechanism that makes a
classification-based argument inductive rather than merely a list of
numerical coincidences.

\subsection{\texorpdfstring{Definition of \(\IMK\)}{Definition of iMK}}

Fix \(S\in\Syl_\ell(X)\) and put
\(\Gamma=\Aut(X)_S\).  The condition \(\IMK(X,\ell)\) asserts that there
are:

\begin{enumerate}
  \item a \(\Gamma\)-stable subgroup
  \[
    \Norm_X(S)\leq N\leq X,
  \]
  with \(N\neq X\) whenever \(\Norm_X(S)\neq X\); and
  \item a \(\Gamma\)-equivariant bijection
  \[
    \Omega:\Irr_{\ell'}(X)\longrightarrow\Irr_{\ell'}(N)
  \]
  such that, for every \(\chi\) with \(\Omega(\chi)=\chi'\),
  \[
    (X\rtimes\Gamma_\chi,X,\chi)
      \geq_c
    (N\rtimes\Gamma_\chi,N,\chi').
  \]
\end{enumerate}

For universal covers of nonabelian finite simple groups, this is
equivalent to the inductive McKay condition of
Isaacs--Malle--Navarro \cite{IsaacsMalleNavarro2007}; see the comparison
in \cite[Section 2.B]{CabanesSpath2026}.

\begin{theorem}[Reduction theorem]\label{thm:reduction}
Fix a prime \(\ell\).  If the universal covering group of every
nonabelian finite simple group satisfies \(\IMK\) at \(\ell\), then every
finite group satisfies \(\IMK\) at \(\ell\).
\end{theorem}

The original reduction is due to Isaacs--Malle--Navarro
\cite{IsaacsMalleNavarro2007}.  The exact centrally-isomorphic-triple
form used here is Rossi's Theorem B \cite{Rossi2023}, recalled as
Theorem 2.10 in \cite{CabanesSpath2026}.

\subsection{Inside the reduction theorem}

For completeness, here is the logical mechanism hidden inside
\cref{thm:reduction}.  A clean equality-level version is Sp{\"a}th's
Theorem 3.15 \cite{Spath2018}; Rossi's stronger proof follows the same
minimal-counterexample spine while retaining the central-isomorphism
data.

One first proves a relative form over each character of the center.
Assume that this form fails and choose a counterexample \(A\) minimizing
\(|A/\Z(A)|\).  The reduction proceeds as follows.

\begin{enumerate}
  \item Local character correspondences and the Glauberman
  correspondence show that both the largest normal \(\ell\)-subgroup
  and the largest normal \(\ell'\)-subgroup of \(A\) lie in \(\Z(A)\).
  Otherwise one passes to a proper quotient or normalizer and contradicts
  minimality.
  \item The generalized Fitting subgroup
  \(F^*(A)=F(A)E(A)\) now has no noncentral solvable part.  A noncentral
  minimal normal section therefore comes from the layer \(E(A)\).  There
  is a normal subgroup \(L\) with
  \[
    L/\Z(A)\cong T^m
  \]
  for a nonabelian finite simple group \(T\) whose order is divisible by
  \(\ell\).
  \item The assumed \(\IMK\) correspondence for the universal cover of
  \(T\) is extended to the \(m\) components, including their permutation
  action, and then descended through the relevant central quotient.
  This produces a proper local subgroup of \(L\), containing the needed
  Sylow normalizer, and a character correspondence compatible with the
  action of \(A\).
  \item Clifford theory lifts that correspondence from \(L\) to \(A\).
  The resulting intermediate subgroup is proper, so minimality applies
  again and replaces it by the actual Sylow normalizer.
  \item Restriction and Butterfly transport align the automorphism
  actions on the two stages; transitivity of central isomorphism composes
  their character-triple witnesses.  This contradicts the choice of
  \(A\).
\end{enumerate}

Thus a hypothetical counterexample would contain a simple section whose
universal cover violates \(\IMK\).  The hypothesis of
\cref{thm:reduction} excludes every such section, proving the result.

\subsection{\texorpdfstring{Why \(\IMK\) gives the actual Sylow normalizer}
  {Why iMK gives the actual Sylow normalizer}}

\begin{proposition}\label{prop:imk-to-mckay}
If every finite group satisfies \(\IMK\) at \(\ell\), then
\cref{thm:mckay} holds at \(\ell\).
\end{proposition}

\begin{proof}
We prove a stronger statement by induction on
\[
  m(X,S):=[X:\Norm_X(S)].
\]
If \(m(X,S)=1\), then \(S\trianglelefteq X\) and
\(\Norm_X(S)=X\), so the identity map is the required bijection.

Assume \(m(X,S)>1\).  The condition \(\IMK(X,\ell)\) supplies a subgroup
\[
  \Norm_X(S)\leq N<X
\]
and a bijection
\[
  \Omega_X:\Irr_{\ell'}(X)\xrightarrow{\sim}\Irr_{\ell'}(N).
\]
The subgroup \(S\) is also a Sylow \(\ell\)-subgroup of \(N\): its order
already contains the full \(\ell\)-part of \(|X|\), hence of \(|N|\).
Moreover
\[
  \Norm_N(S)=N\cap\Norm_X(S)=\Norm_X(S),
\]
because \(\Norm_X(S)\leq N\).  Since \(N<X\),
\[
  [N:\Norm_N(S)]<[X:\Norm_X(S)].
\]
The induction hypothesis applied to \(N\) therefore gives a bijection
\[
  \Omega_N:\Irr_{\ell'}(N)
    \xrightarrow{\sim}
  \Irr_{\ell'}(\Norm_N(S))
  =
  \Irr_{\ell'}(\Norm_X(S)).
\]
The composite \(\Omega_N\circ\Omega_X\) proves the desired equality.

For the stronger equivariant statement, one restricts the second outer
action to the image of \(\Gamma\), applies the Butterfly property, and
uses transitivity of \(\geq_c\).  This is the final induction in the proof
of Theorem B of \cite{CabanesSpath2026}.
\end{proof}

\section{Verification on universal covers}

By \cref{thm:reduction,prop:imk-to-mckay}, it remains to verify
\(\IMK\) for the universal covering groups of the nonabelian finite
simple groups.  The classification divides these into alternating
groups, sporadic groups, and groups of Lie type.  The proof is cumulative:
the 2026 paper closes the only family left by the preceding results.

\begin{longtable}{p{0.24\textwidth}p{0.50\textwidth}p{0.18\textwidth}}
\toprule
Family or prime & Result supplying \(\IMK\) & Source \\
\midrule
\endfirsthead
\toprule
Family or prime & Result supplying \(\IMK\) & Source \\
\midrule
\endhead
Alternating and sporadic groups
  & Inductive McKay for simple groups not of Lie type
  & \cite{Malle2008} \\
Groups of Lie type in defining characteristic
  & Inductive McKay in defining characteristic
  & \cite{Spath2012} \\
\(\ell=2\), all families
  & Equivariant character correspondence for odd-degree characters
  & \cite{MalleSpath2016} \\
\(\ell=3\), all families
  & Type-\(D\) extension results, together with the earlier families
  & \cite{Spath2025II} \\
Types \(A\) and \({}^{2}A\)
  & Equivariant character correspondences
  & \cite{CabanesSpath2017A} \\
Type \(C\)
  & Verification of the inductive condition
  & \cite{CabanesSpath2017C} \\
Types \(B,E_6,{}^{2}E_6,E_7\)
  & Descent equalities and local extension conditions
  & \cite{CabanesSpath2019} \\
Exceptional types \({}^{3}D_4,E_8,F_4,{}^{2}F_4,G_2\)
  & Equivariance and extendibility in connected-center groups
  & \cite{CabanesSpath2013} \\
Suzuki and Ree types \({}^{2}B_2,{}^{2}G_2\)
  & Cases contained in the original reduction analysis
  & \cite{IsaacsMalleNavarro2007} \\
Type \(D_l\) and \({}^{2}D_l\), remaining non-defining primes
  & Final local conditions \(\Ad\) and \(\Bd\)
  & \cite{CabanesSpath2026} \\
\bottomrule
\end{longtable}

This inventory is formalized as Theorem 2.11 of
\cite{CabanesSpath2026}: before its new type-\(D\) argument, only
\(D_l\) and \({}^{2}D_l\) at primes \(\ell\geq5\) different from the
defining characteristic remain.

\section{The local criterion for groups of Lie type}

\subsection{Finite reductive setup}

Let \(\G\) be a simple simply connected linear algebraic group over
\(\overline{\mathbb F}_p\), let \(F:\G\to\G\) be a Frobenius
endomorphism defining an \(\mathbb F_q\)-structure, and set
\[
  G=\G^F.
\]
A regular embedding \(\G\leq\widetilde{\G}\) with connected center
produces
\[
  G\trianglelefteq\widetilde G=\widetilde{\G}^{\,F}.
\]
Field and graph automorphisms form a group \(E\), and
\(\widetilde G\rtimes E\) induces the relevant outer automorphisms of
\(G\).

For an odd prime \(\ell\nmid q\), let
\[
  d=d_\ell(q)
\]
be the multiplicative order of \(q\) modulo \(\ell\).  A Sylow
\(d\)-torus \(\Sbf\) of \((\G,F)\) is the finite-reductive analogue of
a Sylow subgroup.  Generic Sylow theory
\cite{BroueMalle1992,MalleTesterman2011} relates
\(\Norm_G(\Sbf)\) to a Sylow \(\ell\)-subgroup of \(G\).  Set
\[
\begin{aligned}
  N&=\Norm_G(\Sbf),\\
  \widehat N&=\Norm_{G\rtimes E}(\Sbf),\\
  \widetilde N&=\Norm_{\widetilde G}(\Sbf),\\
  \widetilde C&=\Cent_{\widetilde G}(\Sbf).
\end{aligned}
\]

\subsection{The three extension conditions}

The global condition \(\Ainf\) says that there is an \(E\)-stable
\(\widetilde G\)-transversal in \(\Irr(G)\) whose members extend to
their stabilizers in \(G\rtimes E\).  Sp{\"a}th's Theorem A
\cite{Spath2025II}, together with the earlier type-\(A\), type-\(C\), and
type-\(B/E\) results, establishes \(\Ainf\) in all types;
\cite{Spath2023I,Spath2025II} supplies the final type-\(D\) part.

The local condition \(\Ad\) says that there is a
\(\widehat N\)-stable \(\widetilde N\)-transversal in \(\Irr(N)\)
whose members extend to the appropriate stabilizers.  The local
condition \(\Bd\) says:

\begin{enumerate}
  \item maximal extendibility holds for
  \(N\trianglelefteq\widetilde N\) and
  \(\widetilde C\trianglelefteq\widetilde N\); and
  \item there is an extension map for
  \(\widetilde C\trianglelefteq\widetilde N\) equivariant under both
  the linear-character twists of \(\widetilde G/G\) and the outer
  normalizer action.
\end{enumerate}

The local criterion, Theorem 2.21 of \cite{CabanesSpath2026} (imported
from \cite{CabanesSpath2019}), uses the previously established
\(\Ainf\) theorem through Theorem 2.18.  In dependency-explicit form it
states:

\begin{proposition}[Local criterion]\label{prop:local-criterion}
Assume \(\ell\nmid 6q\), \(G/\Z(G)\) is simple, \(G\) is its universal
cover, \(\Ainf\) has been established, and \(\Ad\) and \(\Bd\) hold for
\(d=d_\ell(q)\).  Then \(G\) satisfies \(\IMK\) at \(\ell\), with
local subgroup \(\Norm_{\G}(\Sbf)^F\).
\end{proposition}

The proof combines \(\Ainf\), the local extension data, and a common
parametrization of global and local characters obtained through
generic Harish--Chandra theory and Jordan decomposition
\cite{Malle2007,BroueMalleMichel1993}.  The resulting bijection is
outer-equivariant and its paired characters define centrally
isomorphic character triples.

\section{\texorpdfstring{The final type-\(D\) verification}
  {The final type-D verification}}

\subsection{Arithmetic organization}

Take
\[
  G=D_{l,\mathrm{sc}}^\varepsilon(q),
  \qquad l\geq4,\quad \varepsilon\in\{1,-1\}.
\]
The polynomial order of \((\G,F)\) is
\[
 P_{(\G,F)}(X)
 =
 X^{l^2-l}
 (X^2-1)(X^4-1)\cdots(X^{2l-2}-1)(X^l-\varepsilon).
\]
Let \(a_{(\G,F)}(m)\) be the multiplicity of a primitive \(m\)-th root
of unity as a root of this polynomial.

An integer \(d\geq3\) is \emph{doubly regular} for \((\G,F)\) when
\[
  d\mid 2l
  \quad\text{and}\quad
  \frac{2l}{d}
  \begin{cases}
    \text{is even},&\varepsilon=1,\\
    \text{is odd},&\varepsilon=-1.
  \end{cases}
\]
This is stronger than ordinary regularity: it controls the centralizer
of a Sylow \(\ell\)-subgroup both in the type-\(D_l\) group and in the
ambient type-\(B_l\) overgroup.

\subsection{Doubly regular case}

Theorem 4.1 of \cite{CabanesSpath2026} proves that if \(d\) is doubly
regular, then \(\Ad\) and \(\Bd\) hold.  Its construction uses:

\begin{enumerate}
  \item explicit Tits-group representatives in the torus normalizer;
  \item relative Weyl groups
  \(\Norm_{\G}(\Sbf,\xi)^F/\Cent_{\G}(\Sbf)^F\);
  \item compatible extension maps for characters of the torus
  centralizer; and
  \item control of field, graph, and diagonal automorphisms.
\end{enumerate}

Consequently \cref{prop:local-criterion} gives \(\IMK\) for every odd
prime \(\ell\nmid q\) with \(d_\ell(q)=d\).  Since \(d\geq3\), such a
prime is at least \(5\).  The same theorem separately settles rank
\(l=4\).

\subsection{The trichotomy}

For \(l\geq5\) and \(d\geq3\), Lemma 5.8 of
\cite{CabanesSpath2026} proves that at least one of the following three
alternatives holds:

\begin{enumerate}[label=(\alph*)]
  \item \(d\) is doubly regular;
  \item \(a_{(\G,F)}(d)\leq1\), and
  \(a_{(\G,F)}(dm)=0\) for every odd \(m\geq3\); or
  \item there are \(l_1,l_2>0\), signs
  \(\varepsilon_1,\varepsilon_2\), and \(j\in\{1,2\}\) such that
  \[
    l_1+l_2=l,\qquad
    \varepsilon_1\varepsilon_2=\varepsilon,
  \]
  the \(j\)-th type-\(D\) factor has rank \(l_j\geq4\), \(d\) is
  doubly regular for that factor, and its \(d\)-multiplicity equals
  \(a_{(\G,F)}(d)\).
\end{enumerate}

Case (a) is already done.  For case (b), the needed cyclotomic step is
the following: for odd \(\ell\),
\[
  \ell\mid\Phi_m(q)
  \quad\Longleftrightarrow\quad
  m=d\ell^a\text{ for some }a\geq0
\]
\cite[Lemma 5.2(a)]{Malle2007}.  Alternative (b), applied with
\(m=\ell^a\), gives
\(a_{(\G,F)}(d\ell^a)=0\) for \(a\geq1\).  Thus the entire
\(\ell\)-part of \(|G|\) lies in a Sylow \(d\)-torus of rank at most one,
so the Sylow \(\ell\)-subgroup is cyclic.  The inductive Alperin--McKay
theorem for cyclic-defect blocks \cite{KoshitaniSpath2016} implies the
required \(\IMK\) condition.

\subsection{\texorpdfstring{The subgroup \(\M\) in the non-doubly-regular case}
  {The subgroup M in the non-doubly-regular case}}

In case (c), Cabanes and Sp{\"a}th pass to a twisted Frobenius
\[
  F'=\nu F_q
\]
and construct a generally disconnected maximal-rank algebraic subgroup
\(\M\leq\G\), stable under \(F'\), such that
\[
  (\M^\circ)_{\mathrm{der}}
  \text{ has type }
  D_{l_1}\times D_{l_2},
  \qquad
  |\M/\M^\circ|=2.
\]
At fixed points,
\[
  M_0=G_1.G_2\trianglelefteq
  M^\circ=(\M^\circ)^{F'}
\]
is a central product.  Only the distinguished factor \(G_j\) is
guaranteed to have rank \(l_j\geq4\); the other can have low rank and
need not be quasisimple.  The distinguished factor contains a Sylow
\(d\)-torus with the same \(d\)-multiplicity as \(G\), and \(d\) is
doubly regular there.  Thus its already proved correspondence can be
transported into \(M=\M^{F'}\).

The technical core is Clifford theory across
\[
  M_0=G_1.G_2
   \trianglelefteq M^\circ
   \trianglelefteq M,
\]
where
\[
  [M:M_0]=2\gcd(2,q-1);
\]
the quotient therefore has order \(4\) for odd \(q\), but order \(2\)
for even \(q\).  Characters are partitioned according to their
stabilizers and extension behavior.  Theorem 5.20 of
\cite{CabanesSpath2026} constructs an \(E(M)\)-stable
\(\widetilde M\)-transversal in \(\Irr(M)\), and its chosen characters
extend to the required stabilizers.  In the exceptional case
\(2\mid f\), \(\varepsilon=1\), and \(\varepsilon_1=-1\), the theorem
also controls the central element \(h_0\): when \(h_0\in\ker(\chi)\), it
chooses an extension with \(vF_q\) in its kernel.

Proposition 6.3 transports the needed character correspondences using
central isomorphisms of character triples.  Theorem 6.4 proves
extendibility for the chosen transversal with respect to
\(N\trianglelefteq\widehat N\), and Corollary 6.5 proves maximal
extendibility for \(N\trianglelefteq\widetilde N\).  Propositions 6.7
and 6.8 construct the required equivariant extension maps.  Theorem 6.9
assembles these results for the normalizer of a Sylow \(d\)-torus in
\(G\), proving:

\begin{enumerate}
  \item a \(\widehat N\)-stable \(\widetilde N\)-transversal of
  \(\Irr(N)\), whose members extend to their stabilizers in
  \(\widehat N\);
  \item maximal extendibility for \(N\trianglelefteq\widetilde N\);
  \item on a chosen \(\widehat N\)-stable
  \(\widetilde C\)-transversal in \(\Irr(C)\), a
  \(\widehat N\)-equivariant extension map for
  \(C\trianglelefteq N\); and
  \item a
  \(\operatorname{Lin}(\widetilde N/N)\rtimes\widehat N\)-equivariant
  extension map for
  \(\widetilde C\trianglelefteq\widetilde N\), which also gives maximal
  extendibility for that inclusion.
\end{enumerate}

The first item is \(\Ad\).  The second and fourth items give \(\Bd\);
the third is the auxiliary extension map used in their construction.
Therefore \cref{prop:local-criterion} applies in every
non-doubly-regular case with \(a_{(\G,F)}(d)\geq2\).

\begin{proposition}[Final type-\(D\) case]\label{prop:type-d}
Let \(q\) be a prime power and let \(\ell\nmid 2q\) be prime.  Then
\(\IMK\) holds for \(D_{l,\mathrm{sc}}^\varepsilon(q)\), for every
\(l\geq4\) and \(\varepsilon\in\{1,-1\}\).
\end{proposition}

\begin{proof}
For \(l=4\), use the rank-four clause of the doubly regular theorem.
For \(l\geq5\), put \(d=d_\ell(q)\).  The cases \(d=1,2\) were already
proved in the odd-degree/local-extension work
\cite{MalleSpath2016}.  Thus assume \(d\geq3\) and apply the trichotomy.
The doubly regular case follows from Theorem 4.1; the multiplicity-at-most
one case has cyclic Sylow subgroup and follows from
\cite{KoshitaniSpath2016}; the remaining case follows from Theorem 6.9
and the local criterion.  These cases are exhaustive.
\end{proof}

This is Theorem 6.12 of \cite{CabanesSpath2026}.  Lemma 6.10 chooses the
ranks, signs, and subgroup \(M\), and proves the Sylow-\(d\)-torus and
normalizer containments in the twisted setup.  Lemma 6.11 constructs the
conjugation/isomorphism transport back to the original \(F\)-fixed group.

\section{Assembly of the full theorem}

\begin{proof}[Proof of \cref{thm:mckay}]
Fix a prime \(\ell\).  By the classification of finite simple groups,
every nonabelian finite simple group is alternating, sporadic, or of Lie
type.  The results in the table of Section 3 verify \(\IMK\) for the
universal covers of all families except the remaining type-\(D\) cases.
Those cases are supplied by \cref{prop:type-d}.  Hence the universal
cover of every nonabelian finite simple group satisfies \(\IMK\) at
\(\ell\).

Apply the reduction theorem, \cref{thm:reduction}.  Every finite group
now satisfies \(\IMK\) at \(\ell\).  Apply
\cref{prop:imk-to-mckay} to obtain a bijection between
\(\Irr_{\ell'}(X)\) and
\(\Irr_{\ell'}(\Norm_X(S))\) for every finite \(X\) and
\(S\in\Syl_\ell(X)\).  Taking cardinalities proves the equality.
\end{proof}

\section{Was a simpler proof available?}

We searched the post-2024 literature and the cited character-correspondence
literature for a proof that avoids the inductive condition or the
classification.  No such proof of the theorem for \emph{all} finite
groups was found.  The Annals authors themselves organize the result as
the last step of a classification-based program
\cite{CabanesSpath2026}.

There are important simpler proofs in restricted settings:

\begin{itemize}
  \item the Okuyama--Wajima argument proves the blockwise height-zero
  counting equality, hence McKay's equality, for \(p\)-solvable groups
  \cite{OkuyamaWajima1980};
  \item the prime \(2\) admits a much more focused Harish--Chandra
  analysis \cite{MalleSpath2016}; and
  \item cyclic Sylow/defect cases follow from the cyclic-defect theory
  \cite{KoshitaniSpath2016}.
\end{itemize}

None of these covers the non-\(p\)-solvable type-\(D\) families that
constituted the final obstruction.  Treating a special-case proof as a
general proof would therefore leave a genuine logical gap.  The proof
above is the complete route documented by the cited primary literature:
it compresses the classification work behind named theorems, but does
not disguise or omit it.

There is nevertheless a shorter route for one important part of a
\emph{numerical} formalization.  Using Gallagher's original
commutator-sum method in the explicit form described by Murray, together
with the classical ordinary statement and Okuyama--Wajima comparison
\cite{Gallagher1971,Murray2023,Navarro2018,OkuyamaWajima1980}, good-class counting can
be applied inside each actual character inertia group.  This direct route
uses Clifford multiplicities, Schur orthogonality, and finite Fourier
orthogonality, and avoids constructing a strong character-triple
isomorphism solely for the counting step.  For
\(\theta\in\Irr_{p'}(CS)\), let \(r\) be its canonical invariant
\(p'\)-kernel constituent.  The two groups
\(\Delta=I_X(r)\) and \(\Gamma_\theta=I_X(\theta)\) must not be
conflated: \(\theta\) need not be \(\Delta\)-invariant, and the
per-character count is applied in \(\Gamma_\theta\).  There the relevant
action is self-conjugation, so no external fixed-character theorem is
needed.  Normal-product
restriction and a Schur--Zassenhaus complement first compare the fibres
above the two canonical base extensions.  The source and target Gallagher
linear parameters can then be removed by two independent twist
equivalences.  Only existence of each ambient linear extension is needed;
the two parameters need not be identified with one another.

The Okuyama--Wajima input is used in its printed normal-product scope.
The extension comparison for arbitrary intermediate subgroups is derived
prime-locally: Navarro's Sylow-preimage criterion handles localization and
reassembly, while the original theorem is applied only on the
prime-away local groups where its Sylow, normalizer, intersection, and
abelian-quotient hypotheses hold.  The projective multiplicity
calculation has a second important bookkeeping boundary.  Its retained
prime-to-\(p\) factor is defined on the acting group, whereas the character
extension obstruction lives on the common outer quotient.  Equality of the
two quotient \(H^2\)-classes does construct the required literal factor
matching by a cochain gauge, but vanishing after inflation does not by
itself imply quotient vanishing.  This residue route would therefore
require a separate descent theorem, and no unconditional injectivity of
inflation is assumed.

The printed proof does not establish that stronger residue statement.
It proves extension equivalence directly by induction on the ambient group
order.  After choosing a \(p\)-complement \(M\) in the Sylow normalizer, a
Gorenstein subgroup theorem
\cite[Chapter~3, Theorem~2.2]{Gorenstein1968} supplies
\(L\leq U\trianglelefteq M\) with \(M/U\) cyclic and
\(C_{P/P'}(U)\neq1\).  If the fixed subgroup is all of \(P/P'\), apply
Isaacs Theorem~13.29
\cite[Theorem~13.29]{Isaacs1976} to the action-invariant part of the
source lying-over fibre.  It identifies that subfibre with the complete
fixed-point target fibre.  A target extension and the abelian quotient
give the target fibre cardinality equal to the subgroup index.  The
invariant source subfibre injects into the complete source fibre, while
Clifford theory bounds the latter above by the same index.  Equality
follows, and the sum-of-squared-restriction-multiplicities identity forces
every source constituent to be an extension.  This avoids any need to
strengthen Theorem~13.28 from existence of one invariant constituent to
invariance of the whole source fibre.  Cyclic extension and
prime-power-index ascent then give the conclusion.  Otherwise one forms the
proper nontrivial subgroup \(Q/P'=C_{P/P'}(U)\), uses Glauberman
transitivity, and applies induction to \(N_G(Q)/Q\) and to \(QMK\), with
the same prime-power-index ascent between them.  This direct induction is
the source-faithful route implemented in the formalization; the residue
construction remains an optional sufficient-condition adapter.

This correction matters because the tempting equivariant theorem is not a
valid dependency.  Maltempo--Vallejo Theorem~1.3
\cite{MaltempoVallejo2026} is false as printed: take
\(A=S_3\), \(G=A_3\), \(N=1\), and \(\theta=1\).  Every element of \(A_3\)
is \(\theta\)-good, but only the trivial irreducible character of \(A_3\)
is \(S_3\)-fixed.  The published proof incorrectly asserts that every
restriction from \(S_3\) lies in the span of individually fixed
irreducibles; the degree-two character restricts as
\(\omega+\omega^2\).  Consequently its Theorem~1.3, Corollary~1.4, and
the downstream equivariant counting statements are not used here.

The ordinary count is followed by the now-formalized relative-degree
theorem \cite[Theorem~5.12]{Navarro2018} before passing to the project's
\(p'\)-degree fibre types.  Applied inside \(I_G(\theta)\), that theorem
states that \(\chi(1)/\theta(1)\) divides
\(|I_G(\theta):N|\); the relevant \(p'\)-order input is therefore the
inertia quotient, not the distinct ambient induction index
\(|G:I_G(\theta)|\).  The formalization also identifies the ambient normal
\(p'\)-kernel with its internal copy in the existing Theorem~4.4 API, and
uses equivariance plus injectivity to identify the target inertia group
with the normalizer of the induced Sylow subgroup inside the source
inertia group by a canonical equivalence between the corresponding
subtype groups.
After adjoining the ambient center, arbitrary equivalences between the
resulting finite fibres automatically preserve the central scalar needed
by the numerical reduction.  Thus equal finite cardinalities can be
converted to equivalences without formalizing the
Dade-algebra/projective-obstruction construction.

This shortcut does not replace the classification or the quasisimple
verification.  The efficient formal route is hybrid: retain the
CFSG/\(\IMK\) layer reduction, but use the valid ordinary
Gallagher--Okuyama--Wajima argument for the solvable normal-subgroup seam.
That counting seam is now closed; the remaining mathematical
formalization is the quasisimple/CFSG verification.

\section{Dependency summary}

The logical spine of the argument is:
\[
\begin{gathered}
\text{ordinary character theory and character triples}\\
\Downarrow\\
\text{family-specific equivariant bijections and extension maps}\\
\Downarrow\\
\IMK\text{ for every universal cover (CFSG case split)}\\
\Downarrow\quad\text{\cite{IsaacsMalleNavarro2007,Rossi2023}}\\
\IMK\text{ for every finite group}\\
\Downarrow\quad\text{induction on }[X:\Norm_X(S)]\\
|\Irr_{\ell'}(X)|
=
|\Irr_{\ell'}(\Norm_X(S))|.
\end{gathered}
\]

\begingroup
\footnotesize
\bibliographystyle{alpha}
\bibliography{references}
\endgroup

\end{document}
